\documentclass{article}
\usepackage[utf8]{inputenc}
\usepackage{amsmath}
\usepackage{amssymb}
\usepackage{geometry}
\geometry{a4paper, margin=1in}

\title{核心可解数学模型}
\author{Paper 8}
\date{\today}

\usepackage[UTF8]{ctex}
\begin{document}

\maketitle

\section*{保留的核心模型}
根据\texttt{【tex}公式以及PDF文件，可以直接求解的核心模型是紧凑型MILP重构形式$R$-MILP，适用于情况$M=2$，具体如定理3（第14页和第18页）所述。早期的B-IFP和R-BFP公式被省略了，因为它们属于分数/非线性中间模型；而后期的外部近似算法则属于求解方法，而非核心模型本身。

\paragraph{决策变量。}
$x_j \in \{0,1\}$开启无人机基地$j$，$y_{ij}\in\{0,1\}$将过量请求位置$i$分配给基地$j$，而$\gamma_j^{(m)}\in\{0,1\}$表示是否恰好有$m$架无人机位于基地$j$，用于$m=1,2$。其余变量$U_j^{(1)},U_j^{(2)},z_{jlt},\mu_{lj}^{(m)},\tau_{jlt},\omega_{ij}^{(m)}$是由精确线性化引入的标准辅助变量。设$\bar U_j^{(m)}$是$U_j^{(m)}$的有效上界，且$\bar\mu_{ij}^{(1)}$是$\mu_{ij}^{(1)}$的有效上界。

\paragraph{基础可行集合。}
二进制可行集$\mathcal B$是由分配约束和资源限制所定义的。
\begin{align}
\sum_{i \in I_j} \lambda_i y_{ij}\mathbb{E}[S_{ij}] &\le \sum_{m=1}^{2} m\gamma_j^{(m)} - \epsilon
&& \forall j \in J \tag{9b}\\
x_j \le \sum_{m=1}^{2} m\gamma_j^{(m)} &\le 2x_j
&& \forall j \in J \tag{9c}\\
\sum_{j \in J} \sum_{m=1}^{2} m\gamma_j^{(m)} &= p \tag{9d}\\
\sum_{m=1}^{2} \gamma_j^{(m)} &\le 1
&& \forall j \in J \tag{9e}\\
\sum_{j \in J_i} y_{ij} &= 1
&& \forall i \in I \tag{8c}\\
y_{ij} &\le x_j
&& \forall i \in I,\ \forall j \in J_i \tag{8d}\\
\sum_{j \in J} x_j &\le q \tag{8e}\\
x_j,\ y_{ij},\ \gamma_j^{(m)} &\in \{0,1\}
&& \forall i \in I,\ \forall j \in J_i,\ \forall j \in J,\ m=1,2 \tag{8h/9f}
\end{align}

\section*{紧凑化的 MILP 重新表述形式：$M=2$}
\begin{equation}
\min
\sum_{i \in I} \sum_{j \in J_i} \frac{y_{ij} d_{ij} \lambda_i}{v \sum_{l \in I} \lambda_l}
+
\sum_{i \in I} \sum_{j \in J_i}
\left[
\frac{\omega_{ij}^{(1)}}{2} + \frac{\omega_{ij}^{(2)}}{2}
\right]
\frac{\lambda_i}{\sum_{l \in I} \lambda_l}
\tag{13a}
\end{equation}

\begin{equation}
(x,y,\gamma) \in \mathcal B
\tag{13b}
\end{equation}

\begin{equation}
U_j^{(1)} =
\sum_{l \in I_j} \lambda_l \mu_{lj}^{(1)} \tilde S_{lj}
+
\sum_{l \in I_j} \lambda_l y_{lj} \tilde S_{lj}^{2}
\quad \forall j \in J
\tag{13c}
\end{equation}

\begin{equation}
4U_j^{(2)} =
\sum_{l \in I_j} \lambda_l^2 \tilde S_{lj}
\left(\mu_{lj}^{(2)} \tilde S_{lj} + y_{lj}\tilde S_{lj}^{2}\right)
+
2\sum_{(l,t)\in D_j}\lambda_l\lambda_t \tilde S_{tj}
\left(z_{jlt}\tilde S_{lj}^{2} + \tau_{jlt}\tilde S_{lj}\right)
\quad \forall j \in J
\tag{13d}
\end{equation}

\begin{equation}
(y_{lj}, y_{tj}, z_{jlt}) \in \mathcal F_y^z
\quad \forall (l,t)\in D_j,\ \forall j \in J
\tag{13e}
\end{equation}

\begin{equation}
(y_{lj}, U_j^{(m)}, \mu_{lj}^{(m)}) \in \mathcal M_{yU}^{\mu}
\quad \forall l \in I_j,\ \forall j \in J,\ m=1,2
\tag{13f}
\end{equation}

\begin{equation}
(z_{jlt}, U_j^{(2)}, \tau_{jlt}) \in \mathcal M_{zU}^{\tau}
\quad \forall (l,t)\in D_j,\ \forall j \in J
\tag{13g}
\end{equation}

\begin{equation}
(\gamma_j^{(1)}, \mu_{ij}^{(1)}, \omega_{ij}^{(1)}) \in {\mathcal M'}_{\gamma\mu}^{\omega}
\quad \forall i \in I,\ \forall j \in J_i
\tag{13h}
\end{equation}

\begin{equation}
\omega_{ij}^{(2)} = \mu_{ij}^{(2)}
\quad \forall i \in I,\ \forall j \in J_i
\tag{13i}
\end{equation}

\section*{扩展的线性化集合}
\begin{equation}
z_{jlt} \ge 0,\qquad
z_{jlt} \ge y_{lj} + y_{tj} - 1,\qquad
z_{jlt} \le y_{lj},\qquad
z_{jlt} \le y_{tj}
\quad \forall (l,t)\in D_j,\ \forall j \in J
\tag{Lin-F}
\end{equation}

\begin{equation}
\mu_{lj}^{(m)} \ge 0,\qquad
\mu_{lj}^{(m)} \ge \bar U_j^{(m)}(y_{lj}-1) + U_j^{(m)},\qquad
\mu_{lj}^{(m)} \le \bar U_j^{(m)} y_{lj},\qquad
\mu_{lj}^{(m)} \le U_j^{(m)}
\quad \forall l \in I_j,\ \forall j \in J,\ m=1,2
\tag{Lin-mu}
\end{equation}

\begin{equation}
\tau_{jlt} \ge 0,\qquad
\tau_{jlt} \ge \bar U_j^{(2)}(z_{jlt}-1) + U_j^{(2)},\qquad
\tau_{jlt} \le \bar U_j^{(2)} z_{jlt},\qquad
\tau_{jlt} \le U_j^{(2)}
\quad \forall (l,t)\in D_j,\ \forall j \in J
\tag{Lin-tau}
\end{equation}

\begin{equation}
\omega_{ij}^{(1)} \ge 0,\qquad
\omega_{ij}^{(1)} \ge \bar\mu_{ij}^{(1)}(\gamma_j^{(1)}-1) + \mu_{ij}^{(1)},\qquad
\omega_{ij}^{(1)} \le \bar\mu_{ij}^{(1)} \gamma_j^{(1)},\qquad
\omega_{ij}^{(1)} \le \mu_{ij}^{(1)}
\quad \forall i \in I,\ \forall j \in J_i
\tag{Lin-omega}
\end{equation}

\section*{通过模型强化，保留的削减措施得到了加强。}
\begin{equation}
\gamma_j^{(1)} + \gamma_j^{(2)} \ge U_j^{(m)} / \bar U_j^{(m)}
\quad \forall j \in J,\ m=1,2
\tag{20}
\end{equation}

\begin{equation}
U_j^{(1)} \ge U_j^{(2)}
\quad \forall j \in J
\tag{21}
\end{equation}

\begin{equation}
\gamma_j^{(1)} + \gamma_j^{(2)} = x_j
\quad \forall j \in J
\tag{22}
\end{equation}

\end{document}
